\chapter{Literature Review}


\section{LLMs, debate, and persuasive performance}

Recent work indicates that large language models (LLMs) can match or
exceed human performance in online debates. \textcite{salvi2025_chatgpt_debates} report that
when GPT\mbox{-}4 was provided with basic sociodemographic information about its
opponents, it outperformed human debaters in approximately 64\% of
online debates, achieving a 64.4\% higher success rate in producing
greater post-debate agreement \cite{salvi2025_chatgpt_debates}.
When neither human nor AI debaters had access to such personal information,
GPT\mbox{-}4's persuasive performance was statistically indistinguishable from
that of human opponents, indicating that personalization is a key amplifier
of the model's persuasive advantage \cite{salvi2025_chatgpt_debates}.

This line of research complements earlier work on persuasive
dialogue systems and chatbots. \textcite{shi2020_effects_persuasive_dialogues_bots}
investigate persuasive dialogues for charity donations and report that
intelligent conversational agents can significantly influence donation
behaviour, with different bot identities and inquiry strategies producing
systematic differences in persuasive success and perceived traits such as
competence, confidence, warmth, and sincerity
\cite{shi2020_effects_persuasive_dialogues_bots}. Together, these
studies support the view that LLMs and conversational agents are capable
of substantive persuasive impact in interactive, debate-like settings,
rather than merely in one-shot message framing tasks
\cite{salvi2025_chatgpt_debates,shi2020_effects_persuasive_dialogues_bots}.

\section{Measuring attitude and confidence change}

Persuasion studies involving AI systems commonly operationalize persuasiveness as the change
in self-reported attitudes before and after exposure to an argument or conversation.
Large-scale experiments on AI-driven political persuasion measure participants' positions on
policy issues on continuous or Likert-type scales prior to treatment, expose them to messages
generated or curated with the help of language models, and then re-measure attitudes to quantify
within-subject change \cite{argyle2025_testing_political_persuasion}.
\textcite{durmus2024_measuring_persuasiveness}'s framework for "measuring the persuasiveness of
language models" adopts a similar pre--post design: participants state their views on a claim,
read an argument produced by a model or a human, and then provide a new rating; the difference
is interpreted as the model's persuasive effect \cite{durmus2024_measuring_persuasiveness}.

Beyond immediate attitude shifts, recent work examines the durability of attitude change
following conversational interactions with language models.
\textcite{pataranutaporn2025_oceanchat} and related work report that conversational, narrative-rich
interactive agents can produce measurable increases in behavioral intentions and, in
some cases, durable attitude shifts on targeted outcomes. In OceanChat's large
between-subjects experiment, conversational character narratives produced stronger
behavioral intentions than static informational treatments, while effects on deeper
policy attitudes were more limited \cite{pataranutaporn2025_oceanchat}. This pattern aligns with
evidence that vivid, narrative interactions with LLM-powered agents can support
short-to-medium term attitude or intention change even when long-term policy
positions remain stable.
\section{Statistical inference choices for persuasion models}

Most persuasion studies in this domain estimate average treatment effects using
regression-based comparisons of pre--post outcomes and covariates, then assess
robustness under alternative specifications
\cite{argyle2025_testing_political_persuasion,durmus2024_measuring_persuasiveness}.
For debate datasets with moderate sample size and potentially heterogeneous
variance across conditions, linear models remain appropriate for interpretable
effect estimation, while robust standard errors are commonly used to prevent
overconfident inference under heteroskedasticity.

Relatedly, robustness analysis is increasingly treated as a design requirement
rather than a post-hoc add-on: model conclusions are expected to be checked for
sensitivity to influential observations, alternative covariate sets, and
distributional assumptions. This motivates combining a confirmatory model family
with secondary and exploratory sensitivity analyses, so that inference is based
on stability of effect direction and uncertainty intervals rather than on a
single fitted specification. In the present dataset, outcomes are participant-level
continuous change scores, condition cells are modest and mildly imbalanced, and
heteroskedasticity is plausible across debate conditions and topics; OLS therefore
provides interpretable mean-difference estimates while HC3 robust standard errors
reduce small-sample overconfidence in uncertainty intervals.

\section{Source identity, AI disclosure, and trust}

A separate but related body of research investigates how beliefs about the source of a
message---in particular, whether it is human- or AI-generated---shape trust and evaluations of
content. \textcite{schilke2025_transparency_dilemma} document a "transparency dilemma": across
thirteen preregistered studies in domains such as grading, hiring, investment advice, design,
and routine communication, actors who disclosed using generative AI were trusted less than those
who did not, even when the objective quality of the work was held constant
\cite{schilke2025_transparency_dilemma}. The trust penalty is partly explained by reduced
perceptions of legitimacy, as AI-assisted work is seen as deviating from appropriate
professional norms \cite{schilke2025_transparency_dilemma}.

These findings interact with broader work on trust in humans versus artificial agents.
Research on trust toward humans and AI systems suggests that trust in AI is related to, but not
reducible to, dispositional trust in humans, and that people differentiate between human and
machine sources when forming expectations about reliability and responsibility
\cite{montag2023_trust_humans_vs_ai}. Studies of persuasive chatbots likewise show that bot
"identity" and framing (for example, as an expert advisor versus a peer) influence both
persuasion outcomes and interpersonal perceptions such as warmth, competence, and sincerity
\cite{shi2020_effects_persuasive_dialogues_bots}. Collectively, this literature implies that
perceived source identity (human vs.\ AI) and disclosure practices can alter trust, perceived
legitimacy, and openness to influence, even when message content is held constant
\cite{schilke2025_transparency_dilemma,montag2023_trust_humans_vs_ai,shi2020_effects_persuasive_dialogues_bots}.

\section{Conversational style and AI debater personality}

Beyond source identity, a growing literature examines how the conversational style or personality
of a chatbot shapes user perceptions and persuasive impact. Drawing on the warmth-competence
framework from social psychology, \textcite{deng2022_ambivalent_stereotypes_persuasion} and
related work argue that communicators can be located along two primary dimensions (perceived
warmth (friendliness, care) and competence (ability, expertise)) which jointly influence persuasion
by affecting both source evaluations and message processing
\cite{deng2022_ambivalent_stereotypes_persuasion}. In the context of text-based chatbots,
\textcite{schwede2023_can_chatbots_be_persuasive} show that warmth and competence perceptions can
be systematically manipulated through communication style, such as humanizing language and social
cues \cite{schwede2023_can_chatbots_be_persuasive}.

Experimental evidence supports the idea that distinct chatbot styles have differential persuasive
effects. Prior work shows that warm and competent communication styles can outperform neutral
styles in persuasion tasks and can systematically shift perceived source qualities such as
warmth, competence, and trust; lexical and conversational cues are sufficient to produce these
personality impressions in text-based interaction \cite{schwede2023_can_chatbots_be_persuasive,heppner2024_chatbot_personality_cues}.
The present study manipulates personality as a between-subjects factor but does not collect
warmth/competence perception measures; personality effects are assessed through belief and
confidence outcomes rather than perceptual mediators.

\section{Source attribution and within-debate stance dynamics}

While the literature establishes that pre--post stance measurement is standard, an emerging and
underexplored dimension is how participants' beliefs about opponent identity evolve \textit{during} a
debate, and whether real-time stance tracking across debate rounds reveals persuasion dynamics
masked in summary measures. Research on source attribution indicates that people actively infer
source characteristics from message content, and that misattribution—for instance, believing an
AI-generated argument is human-authored or vice versa—shapes both immediate receptivity and
post-hoc rationalization \cite{schilke2025_transparency_dilemma,montag2023_trust_humans_vs_ai}.
Moreover, rounds of back-and-forth debate create dynamic conditions under which participants
may revise their opponent beliefs and simultaneously adjust their stance; tracking both processes
in parallel can illuminate how perceived source identity moderates incremental persuasion.
Per-round measurement of stance, confidence, and evolving opponent perception—combined with
post-debate justification of those inferences—provides a novel data layer for understanding
how AI versus human argumentation shapes persuasion trajectories over an extended interaction.

\section{Comprehension checks and between-round belief measurement}

While per-round stance tracking captures \textit{when} persuasion occurs, complementary measurement
approaches help validate \textit{what} drove that persuasion. Research on argument comprehension
and learning demonstrates that checking whether individuals can accurately reiterate an opponent's
position reveals depth of engagement and argument processing
\cite{sternberg2013_cognitive_psychology_book}. When combined with contemporaneous belief measurement,
comprehension-checking prevents conflating genuine persuasion (actual change in conviction) with
superficial agreement or inattention. In educational contexts, elaborative retrieval tasks—where
learners restate material in their own words—are associated with durability of attitude change
and reduced susceptibility to backsliding \cite{craik_lockhart_1972_levels_of_processing}.

Similarly, research on metacognitive monitoring shows that explicit, prompted reflection on one's
own beliefs during an ongoing interaction strengthens the connection between argument exposure and
attitude change, because reflection amplifies awareness of the reasoning that supports current beliefs
\cite{schommer1990_effects_reflection_critical_thinking}.
Critically, between-round belief measurement—capturing stance shifts relative to specific AI arguments
rather than relying on gross pre-post change—allows detection of which argument types and rhetorical
strategies more effectively persuade, rather than attributing persuasion to opponent source or
background variables alone \cite{durmus2024_measuring_persuasiveness}.

When applied to interactive debate, between-round belief checks paired with argument comprehension
assessment provide a research-grounded mechanism to (1) validate that participants engaged
substantively with opponent arguments, (2) isolate persuasion to specific argument exchanges rather
than systemic factors, and (3) generate fine-grained data on how AI versus human argumentation
strategies differentially prompt belief revision. These measurements do not themselves constitute
persuasion mechanisms; rather, they serve as methodological tools to isolate and quantify the
persuasive impact of the opponent's argumentation.

\section{Reliability and agreement for AI-assisted quality scoring}

When automated scores are used as explanatory features, methodological quality
depends on demonstrating alignment with human judgement. In practice, two
families of metrics serve complementary purposes: rank-order association
(e.g., Spearman correlation) and inter-rater agreement metrics
(e.g., ICC variants). Rank association evaluates whether the model preserves
relative ordering of higher- vs lower-quality texts, while ICC quantifies
consistency under a shared rating scale and allows direct comparison of
human-only reliability versus human-plus-model reliability.

This distinction is important for debate analysis: a model can track ordering
reasonably well but still exhibit scale disagreement, or conversely can show
moderate absolute agreement with weak rank discrimination. Therefore, feature
validation should not be treated as binary "valid/invalid"; instead it should
be reported as metric-specific evidence with uncertainty-aware interpretation.
This framing directly supports differential confidence across feature families
such as reflection quality versus persuader argument quality.

\section{Implications for debate-format studies of LLM persuasiveness}

Synthesizing these streams of research highlights several design-relevant points for debate-format
investigations of LLM persuasiveness. First, randomized trials of GPT\mbox{-}4 in online debates
demonstrate that LLMs can be as persuasive as, and under personalization more persuasive than,
human opponents, establishing debates as a realistic and informative testbed for AI persuasion
\cite{salvi2025_chatgpt_debates}. Second, pre--post measurement of self-reported stance and
confidence is an established method for quantifying persuasive impact in both AI-mediated political
communication and model-centric evaluation studies, and can be extended to interactive debate
 scenarios \cite{argyle2025_testing_political_persuasion,durmus2024_measuring_persuasiveness,pataranutaporn2025_oceanchat}.
Extending this approach to between-round belief measurement enables detection of persuasion
dynamics at finer granularity: tracking which specific argument exchanges correlate with stance
shifts, and validating that observed changes reflect genuine processing of AI arguments rather
than noise or inattention.

Third, evidence on the transparency dilemma and on trust in AI indicates that whether debaters are
believed to be human or AI, and whether AI use is disclosed, can meaningfully alter trust and
perceived legitimacy, potentially moderating attitude change even when arguments are identical
\cite{schilke2025_transparency_dilemma,montag2023_trust_humans_vs_ai}. Finally, findings on
chatbot personality and communication style suggest that systematically varying an AI debater's
style along warmth and competence dimensions---for example, contrasting a firm, highly assertive
debater with a more balanced and open-minded counterpart---is likely to affect perceived source
qualities and, in turn, persuasiveness
\cite{schwede2023_can_chatbots_be_persuasive,heppner2024_chatbot_personality_cues}. Taken
together, existing work provides a strong empirical and theoretical foundation for studying
LLM-mediated persuasion in debate formats that employ between-round belief measurement to
isolate AI persuasive effects, track stance change relative to specific argument exchanges,
and validate genuine argument processing, while manipulating perceived source identity and
conversational style \cite{salvi2025_chatgpt_debates,argyle2025_testing_political_persuasion,durmus2024_measuring_persuasiveness,schilke2025_transparency_dilemma,schwede2023_can_chatbots_be_persuasive}.

\subsection*{Scope boundary and inferential focus}

The literature motivates both AI-to-human and human-to-AI persuasion dynamics,
but inferential emphasis in this thesis is intentionally asymmetric: confirmatory
testing focuses on AI-to-human outcomes, while human-to-AI dynamics are treated
as a deferred extension. This boundary is important for internal coherence: it
prevents overclaiming beyond the model families and validated features used in
the present analysis.

Observed pilot effect magnitude for pooled belief change ($d=0.286$) is directionally
larger than many political-persuasion benchmarks (often around $d\approx0.1$--$0.2$),
but remains substantially less decisive than win-rate style comparisons such as the
64\% GPT-4 debate advantage reported by \textcite{salvi2025_chatgpt_debates}.
At pilot $n$, confidence intervals remain wide and comparative effects are not
identified with precision.

\subsection*{Construct-to-variable mapping used in analysis}

\begin{table}[h]
\centering
\caption{From Literature Constructs to Model Variables}
\footnotesize
\setlength{\tabcolsep}{4pt}
\renewcommand{\arraystretch}{1.15}
\begin{tabular}{@{}L{2.6cm}L{3.3cm}L{2.6cm}L{5.8cm}@{}}
\toprule
Construct & Operational variable & Used in models & Expected analytical role \\
\midrule
Persuasion outcome & belief\_change & M1--M5 (primary) & Core AI-to-human treatment effect \\
Confidence shift & confidence\_change & conf-M1--conf-M3 (secondary) & Adjacent outcome under null belief effect \\
Argument processing depth & p1\_reflection\_quality & M2, M3, M4, M4b, M5 & Engagement/comprehension covariate \\
Argument quality signal & persuader\_arg\_quality & M4b, M4c, M5 & Exploratory explanatory feature \\
Style/content context & topic indicators, stance strength & M4, M5 & Control for contextual heterogeneity \\
Debate structure & n\_rounds, first\_player & M5 & Structural controls for interaction exposure \\
\bottomrule
\end{tabular}
\end{table}

Accordingly, current literature provides a solid conceptual basis for debate-format
persuasion analysis. Methodological confidence differs by evidence tier:
confirmatory AI-to-human treatment models provide primary inference,
secondary outcomes provide qualified evidence, and exploratory mechanism/feature
analyses are hypothesis-generating.